% Einheit 3 von 4 – Normalform (bank/quadratische-funktionen/e3.jsonl, zusammenbau v0.1)
%% TODO Zweigzeile Teil 1: was hier gelernt wird (Satz in Schülersprache aus dem Einheitstitel) – Einheit 3
\einheitenkopf[e3]{Einheit 3 von 4 · Normalform}
\zweigzeile{ab Kl. 9, je nach Buch bis Kl. 10 · P10 oft}
\setcounter{aufgabe}{14}

%% TODO Titel: Ich-kann-Satz fehlt in der Bank (Platzhalter: Kettenname) – Nr. 15
%% TODO Anweisung: ein Satz über der ersten Teilaufgabe fehlt in der Bank – Nr. 15
\begin{aufgabe}{Normalform}
\begin{teile}
\teil $f(x) = x^2 + 4x - 1$ – welche Form? \kreuz{Scheitelpunktform} \kreuz{Normalform}
\end{teile}
%% TODO Darstellung für schwach fehlt in der Bank (quadratische-funktionen-e3-k1-s1-v1; Abschnitt „Für schwache Schüler“)
\swz{$f(x) = x^2 + 3x + 5$ – Schnittpunkt mit der $y$-Achse?}{}
%% TODO Darstellung für schwach fehlt in der Bank (quadratische-funktionen-e3-k1-s1-v2; Abschnitt „Für schwache Schüler“)
\swz{$f(x) = x^2 - 2x + 7$ – Schnittpunkt mit der $y$-Achse?}{}
%% TODO Darstellung für schwach fehlt in der Bank (quadratische-funktionen-e3-k1-s1-v3; Abschnitt „Für schwache Schüler“)
\swz{$f(x) = x^2 + 6x + 4$ – Schnittpunkt mit der $y$-Achse?}{}
%% TODO Darstellung für schwach fehlt in der Bank (quadratische-funktionen-e3-k1-s1-v4; Abschnitt „Für schwache Schüler“)
\swz{$f(x) = x^2 - 5x + 9$ – Schnittpunkt mit der $y$-Achse?}{}
%% TODO Darstellung für schwach fehlt in der Bank (quadratische-funktionen-e3-k1-s1-v5; Abschnitt „Für schwache Schüler“)
\swz{$f(x) = x^2 + x + 8$ – Schnittpunkt mit der $y$-Achse?}{}
\swfrage{Was bleibt gleich, was ändert sich?}
%% TODO Darstellung für schwach fehlt in der Bank (quadratische-funktionen-e3-k1-s2-v1; Abschnitt „Für schwache Schüler“)
\swz{$f(x) = x^2 + 4x - 2$ – $f(-3)$?}{}
%% TODO Darstellung für schwach fehlt in der Bank (quadratische-funktionen-e3-k1-s3-v1; Abschnitt „Für schwache Schüler“)
\swz{$\text{$f(x) = (x - 2)^2 + 5$ – Normalform?}$}{}
%% TODO Darstellung für schwach fehlt in der Bank (quadratische-funktionen-e3-k1-s4-v1; Abschnitt „Für schwache Schüler“)
\swz{$\text{$f(x) = (x + 4)^2 + 1$ – Normalform?}$}{}
%% TODO Darstellung für schwach fehlt in der Bank (quadratische-funktionen-e3-k1-s5-v1; Abschnitt „Für schwache Schüler“)
\swz{$\text{Stimmt $(x + 5)^2 - 8 = x^2 + 10x + 17$? Weise nach oder widerlege.}$}{}
\swz{$f(x) = x^2 - 4x + 1$ – lies den Scheitel am Graphen ab.}{\begin{ksys}[xmin=-1,xmax=5,ymin=-4,ymax=2,ablesen]
\parabel{1}{2}{-3}{f}
\end{ksys}}
\swz{$f(x) = x^2 - 6x + 5$ – lies den Scheitel ab, stelle die Scheitelpunktform auf und prüfe mit $x = 0$.}{\begin{ksys}[xmin=-1,xmax=6,ymin=-5,ymax=1,ablesen]
\parabel{1}{3}{-4}{f}
\end{ksys}}
\begin{teile}
\teil $f(x) = x^2 + 2x - 3$ – wahr oder falsch? \\ $(2|5)$ liegt auf $f$. \kreuz{wahr} \kreuz{falsch} \\ $f$ ist die um $1$ nach links und um $4$ nach unten verschobene Normalparabel. \kreuz{wahr} \kreuz{falsch} \\ Der Schnittpunkt mit der $y$-Achse ist $(-3|0)$. \kreuz{wahr} \kreuz{falsch}
\end{teile}
\swz[3]{Die Abbildung zeigt die Parabel $f(x) = x^2 - 8x + 13$. Gib den Scheitel an und notiere die Scheitelpunktform. \hfill (P10 2025 OS)}{\begin{ksys}[xmin=-1,xmax=7,ymin=-4,ymax=2,ablesen]
\parabel{1}{4}{-3}{f}
\end{ksys}}
\swz[3]{$\text{Weise nach: $(x + 4)^2 - 3 = x^2 + 8x + 13$. \hfill (P10 2017 OS)}$}{}
\end{aufgabe}

%% TODO Titel: Ich-kann-Satz fehlt in der Bank (Platzhalter: Kettenname) – Nr. 16
%% TODO Anweisung: ein Satz über der ersten Teilaufgabe fehlt in der Bank – Nr. 16
\begin{aufgabe}{Normalform – Fehler finden, Begründen}
\swz[3]{$\text{Nina rechnet: \rechnung{(x - 4)^2 + 1 &= x^2 - 16 + 1 \\ &= x^2 - 15} Finde den Fehler und korrigiere.}$}{}
\swfrage{Der Graph von $f(x) = x^2 - 7x + 4$ schneidet die $y$-Achse im Punkt $(0|4)$. Warum?}
\end{aufgabe}

\uebersichtskasten{Normalform: y = x² + px + q. Die Zahl q ist der y-Achsenabschnitt (x = 0 einsetzen). Die Öffnung ist die der Normalparabel. \\ Von der Scheitelpunktform zur Normalform: Klammer mit der binomischen Formel auflösen, zusammenfassen. \\ (x − 3)² + 1 = x² − 6x + 9 + 1 = x² − 6x + 10 \\ Von der Normalform zur Scheitelpunktform: Scheitel am Graphen ablesen und in (x − d)² + e eintragen; Probe: x = 0 muss q ergeben. \\ x² − 6x + 10 mit Scheitel S(3 | 1) → (x − 3)² + 1; Probe: (0 − 3)² + 1 = 10 (wA)}
